\documentclass[12pt]{article}

\usepackage[margin=1in]{geometry}
\usepackage{amsmath}
\usepackage{hyperref}
\usepackage{enumitem}
\usepackage{titlesec}
\usepackage{booktabs}

\hypersetup{
    colorlinks=true,
    linkcolor=blue,
    citecolor=blue,
    urlcolor=blue
}

\titleformat{\section}{\large\bfseries}{\thesection}{1em}{}
\titleformat{\subsection}{\normalsize\bfseries}{\thesubsection}{1em}{}

\title{Colonial Embargoes and the Jefferson--Madison Embargo:\\
A Historical Overview}
\author{}
\date{}

\begin{document}
\maketitle

\section{Colonial Non-Importation and Non-Exportation Agreements}

As Rothbard details in \emph{Conceived in Liberty} (particularly Volume~4,
covering 1774--1784), the American colonies employed economic coercion as a
principal weapon against British imperial policy well before armed conflict
began.

\subsection{The Non-Importation Agreements (1765--1770)}

In response to the Stamp Act and later the Townshend Acts, colonial merchants
organized boycotts of British goods. Boston, New York, and Philadelphia led
these efforts. The agreements were enforced by local committees---sometimes
coercively---and proved remarkably effective at pressuring British merchants to
lobby Parliament for repeal. The Townshend duties (except on tea) were repealed
in 1770.

\subsection{The Continental Association (1774)}

The First Continental Congress adopted the most sweeping embargo: a three-stage
program of non-importation (effective December 1774), non-consumption (March
1775), and non-exportation (September 1775) of goods to and from Britain. Local
``Committees of Inspection'' enforced compliance, often functioning as shadow
governments. Rothbard emphasizes how these committees became instruments of
revolutionary authority, displacing royal governance at the grassroots.

\subsection{Colony-Specific Actions}

Individual colonies had imposed their own embargoes at various points---Virginia's
Association of 1769, South Carolina's non-importation movement, and others.
Rothbard highlights how these efforts revealed tensions between radical and
conservative factions: merchants who profited from British trade often resisted,
while artisans and inland farmers supported the boycotts enthusiastically.

\subsection{Key Lessons the Founders Drew}

\begin{itemize}[nosep]
    \item Economic coercion could be a substitute for military force.
    \item Embargoes could unify disparate colonies around a common cause.
    \item But enforcement was difficult, smuggling was endemic, and the economic
          pain fell unevenly---themes that would recur dramatically in the 1800s.
\end{itemize}

\section{The Jefferson--Madison Embargo (1807--1809) and Its Successors}

Jefferson and Madison, both steeped in the revolutionary experience, believed
that ``peaceable coercion'' through trade restriction could compel the European
powers, Britain and Napoleonic France, to respect American neutral trading
rights. The logic was explicitly drawn from the colonial precedent: if embargoes
had brought the British Empire to terms in the 1760s--70s, surely the
now-larger American economy could do so again.

The \textbf{Embargo Act of December 1807} prohibited virtually all American
ships from sailing to foreign ports. It was the most dramatic peacetime
restriction on commerce in American history.

\section{Effects on Tariff Revenues and Trade Balances}

The fiscal consequences were devastating:

\begin{itemize}
    \item \textbf{Tariff revenue collapsed.} The federal government depended
    almost entirely on customs duties (roughly 90\% of federal revenue).
    American imports fell from approximately \$138~million in 1807 to
    \$57~million in 1808---a decline of nearly 60\%. Customs revenue fell
    correspondingly, from around \$16~million to roughly \$7--8~million.

    \item \textbf{Exports cratered} from about \$108~million in 1807 to
    \$22~million in 1808, a drop of nearly 80\%. The trade balance, which had
    been modestly negative, became almost irrelevant as both sides of the ledger
    shrank.

    \item \textbf{Government expenditures} could not easily be cut to match.
    The Treasury ran deficits and Secretary Gallatin---who had painstakingly
    built surpluses and paid down the national debt---saw his fiscal position
    badly eroded.
\end{itemize}

After the Embargo Act was replaced by the \textbf{Non-Intercourse Act (1809)}
and then \textbf{Macon's Bill No.~2 (1810)}, trade partially recovered but
remained disrupted. The cycle of partial embargoes and trade restrictions from
1807 to 1812 left federal finances in a weakened state.

\section{Discord Among the States}

The embargo created sectional and political fissures that genuinely threatened
the Union.

\subsection{New England Resistance}

The commercial states of New England---Massachusetts, Connecticut, Rhode
Island---were devastated. Their economies depended on maritime trade, shipping,
and the carrying trade. Federalist politicians denounced the embargo as
unconstitutional tyranny. Town meetings passed resolutions condemning it.
Smuggling across the Canadian border (especially through Vermont and along Lake
Champlain) became rampant and essentially open.

\subsection{Nullification in Embryo}

Several New England state legislatures passed resolutions declaring the Embargo
Act unconstitutional and instructing their citizens and officials not to
cooperate with enforcement. Connecticut's Governor Trumbull refused to deploy the
militia to enforce the embargo. This was a direct precursor to the
\textbf{Hartford Convention (1814--15)}, where New England Federalists would
flirt with secession and propose constitutional amendments to limit
Southern/Republican power.

\subsection{Southern and Western Ambivalence}

Agricultural states in the South and West suffered from the inability to export
tobacco, cotton, and grain, but their political leaders (largely Jeffersonian
Republicans) supported the policy on ideological grounds. This created internal
tensions within those states as well.

\subsection{Enforcement and Civil Liberties}

To combat evasion, Congress passed increasingly draconian enforcement acts. The
\textbf{Enforcement Act of January 1809} gave customs collectors sweeping powers
to seize goods on mere suspicion. Jefferson, the great apostle of limited
government, found himself presiding over what many considered the most intrusive
federal police action since the Alien and Sedition Acts. The irony was not lost
on contemporaries.

\section{Public Finances on the Eve of the War of 1812}

By 1811--1812, the cumulative damage was severe:

\begin{itemize}
    \item \textbf{The national debt}, which Gallatin had reduced from about
    \$83~million in 1801 to roughly \$45~million by 1811, was poised to
    explode. The government could not raise revenue through tariffs at levels
    sufficient to fight a war.

    \item \textbf{The First Bank of the United States} was allowed to expire in
    1811 (its charter was not renewed), partly due to Republican ideological
    hostility. This eliminated a critical institution for government borrowing
    and fiscal management precisely when it was most needed.

    \item \textbf{Military preparedness was abysmal.} The Army was small
    ($\sim$6,700~men), the Navy had been shrunk by Jeffersonian economy, and
    there was no financial infrastructure to scale up rapidly.

    \item \textbf{When war was declared in June 1812}, the Treasury immediately
    faced a borrowing crisis. Gallatin's war loan of \$11~million in 1812 was
    only partially subscribed. Subsequent loans in 1813 and 1814 had to be sold
    at steep discounts---some at 80~cents on the dollar. By 1814, the government
    was essentially bankrupt: it defaulted on interest payments on the national
    debt, Treasury notes traded below par, and different regions of the country
    operated with wildly different effective currencies.

    \item The wartime fiscal crisis ultimately led to the chartering of the
    \textbf{Second Bank of the United States (1816)} and the adoption of the
    \textbf{Tariff of 1816}---the first explicitly protectionist tariff---marking
    a decisive break from the Jeffersonian model of minimal government funded by
    low import duties.
\end{itemize}

\section{The Colonial Precedent vs.\ the Jeffersonian Reality}

The crucial difference between the colonial embargoes and the Jefferson--Madison
embargo was this: in the 1770s, Americans were restricting their consumption of
\emph{British} goods while developing domestic alternatives and smuggling from
other sources. The pain fell significantly on British exporters and merchants. In
1807--1809, Jefferson was restricting \emph{American exports and shipping}---the
pain fell overwhelmingly on Americans themselves, particularly the mercantile
communities of the Northeast. Britain and France, while inconvenienced, had
alternative sources of supply and far larger economies.

The colonial experience had taught the Founders that economic coercion could
work, but the analogy was deeply flawed when applied to a situation where America
was restricting its own commerce rather than boycotting an adversary's goods.

\end{document}
